% Перед заполнением:
% - Запиши целевую профессию, уровень и задачи, которые решает роль.
% - Прочитай вакансию в три прохода: обязательные требования; повторяющиеся
%   обязанности, инструменты и результаты; пожелания отдельно от требований.
% - Оставляй только то, что подтверждено опытом, образованием или проверяемым
%   обучением. Видимый пробел лучше ложного утверждения.
% - Адаптация под вакансию = выбор и порядок правдивых фактов. Нельзя менять
%   даты, должности, работодателей, метрики и масштаб.
% - AI может сравнивать и переформулировать, но обязан сохранить каждый факт,
%   число, должность, дату, технологию и масштаб. Не передавай ему закрытые
%   данные работодателя.
%
% Финальная проверка (make check):
% 1. Факты: даты, должности, работодатели, образование, метрики, ссылки.
% 2. Релевантность: убрать пункты, не помогающие этому решению о найме.
% 3. Язык: прочитать вслух, убрать воду и нерасшифрованные аббревиатуры;
%    настоящее время для текущей работы, прошедшее для завершённой.
% 4. Файл: одна-две страницы, текст извлекается по порядку, ссылки
%    открываются, плейсхолдеров не осталось.

\DocumentMetadata{lang=ru,pdfversion=2.0,tagging=on}
\documentclass[10pt,a4paper]{article}
\usepackage[russian]{babel}
\usepackage{cv}

\begin{document}

\cvheader{Родион Мышалов}{SENIOR BACKEND ENGINEER}{Бангкок, Таиланд}{%
  \href{mailto:rodionmyshalov@gmail.com}{rodionmyshalov@gmail.com}\cvsep
  \href{tel:+66982458188}{+66 98 245 8188}\cvsep
  \cvlink{https://www.linkedin.com/in/myshalove}{linkedin.com/in/myshalove}\cvsep
  \cvlink{https://github.com/myshalove}{github.com/myshalove}}

\section*{О себе}
% 3-4 строки: профессия и уровень, подтверждаемый стаж, одна-две сильные
% стороны, результат или масштаб, который тебя выделяет. Без «цели»,
% без soft skills и «нацеленности на результат».
Backend-разработчик с опытом X лет в области. Построил систему и получил проверяемый результат или масштаб. Сильные стороны: технология, технология, практика.

\section*{Навыки}
% 6-12 релевантных терминов, о которых можешь говорить на собеседовании и
% которые подтверждены опытом ниже. Без оценок и общих качеств.
\cvskillwidest{Брокеры сообщений}
\cvskill{Языки}{Go, SQL}
\cvskill{API}{REST + OpenAPI/Swagger, gRPC + Protobuf, OAuth + JWT, GraphQL, Webhooks}
\cvskill{Базы данных}{PostgreSQL, MySQL, MongoDB, Redis}
\cvskill{Брокеры сообщений}{Apache Kafka, RabbitMQ}
\cvskill{Инфраструктура}{Docker, Kubernetes (K8s), Helm, Nginx / Angie, HashiCorp Vault}
\cvskill{CI/CD}{GitLab CI, GitHub Actions}
\cvskill{Мониторинг}{Prometheus, Grafana, Sentry, ELK (Kibana)}
\cvskill{Архитектура}{Монолит, Микросервисы}
\cvskill{Проектирование}{TDD, DDD, Clean Architecture}
\cvskill{Тестирование}{Unit, интеграционное, race detection}

\section*{Опыт работы}
% Обратная хронология. Для каждой роли: организация, официальная должность,
% город или «Удалённо», месяц и год, строка о незнакомой компании,
% два-шесть пунктов по убыванию релевантности и силы.
%
% Пункт = сильный глагол + работа + проверяемый результат, масштаб или
% ограничение. «Участвовал в», если ответственность была общей.
% Сильные глаголы: Автоматизировал, Внедрил, Запустил, Исправил, Мигрировал,
% Наставлял, Оптимизировал, Переработал, Повысил, Построил, Разработал,
% Руководил, Сократил, Спроектировал, Стандартизировал, Унифицировал,
% Устранил, Ускорил.

\cventry{SOHO.LMS}{Сентябрь 2023 -- н.\,в.}{Senior Go Engineer}{Валенсия, Испания (удалённо)}
\cvcontext{Онлайн-платформа для образования.}
\begin{itemize}
  \item Разрабатываю на Go сервис расписания и онлайн-уроков: повторяющиеся занятия, учёт доступности преподавателей, переносы и запуск видеоуроков. p95 загрузки календаря упал с 1,2 с до 120 мс (−90\%), сервис держит до 3\,000 RPS в часы начала занятий.
  \item Построил платёжный контур на вебхуках Stripe и эквайринга: автоматическое открытие доступа к курсу после оплаты, абонементы и учёт задолженностей. Доступ выдаётся за 10 секунд вместо 2--4 часов ручной проверки, доля оплат без доступа снизилась с 2,5\% до 0,1\%.
  \item Разработал сервис уведомлений на Kafka с ретраями и дедупликацией: Telegram, WhatsApp, SMS, email и push в мобильное приложение. До 2 млн сообщений в сутки, доля недоставленных --- с 6\% до 0,5\%, посещаемость уроков выросла на 15\% за счёт напоминаний.
\end{itemize}

\cventry{Taptima}{Декабрь 2020 -- Сентябрь 2023}{Middle+ Go Developer}{Лимассол, Кипр (офис и удалённо)}
\cvcontext{Аутсорс-разработка: веб-сервисы, мобильные приложения и e-commerce.}
\begin{itemize}
  \item Перевёл монолитный бэкенд сети кинотеатров на микросервисы и разработал на Go сервис фоновой синхронизации: он забирает из внешней системы 500+ фильмов и 12\,000+ сеансов по отдельности, связывает их между собой и кэширует в PostgreSQL. p95 выдачи афиши упал с 900 до 70 мс (−92\%), сервис держит до 4\,000 RPS в дни премьер, а алерты в чат разработки сократили время обнаружения сбоя синхронизации с нескольких часов до 2 минут.
  \item Перестроил оплату билетов с синхронного REST на асинхронную обработку через очередь задач и вебхуки банка об успешных и неуспешных платежах: p95 ответа при оформлении заказа сократился с 3,5 с до 220 мс (−94\%), доля списаний без выданного билета снизилась с 1,8\% до 0,1\%, обращения в поддержку по оплате --- на 70\%; до 1\,200 заказов в минуту на пике.
  \item Вынес авторизацию в отдельный сервис, который выпускает JWT независимо от остального бэкенда: вход по одноразовым кодам через SMS и email, OAuth 2.0 через Google, Facebook и Telegram. Сервис выдерживает до 2\,000 RPS на логин в часы старта продаж при p99 90 мс, нагрузка на основной бэкенд снизилась на 35\%, а доля ошибок входа на пике --- с 4\% до 0,2\%.
\end{itemize}

\cventry{DNS}{Октябрь 2019 -- Декабрь 2020}{Junior Go Developer}{Владивосток, Россия (офис)}
\cvcontext{Онлайн-магазин цифровой техники.}
\begin{itemize}
  \item Разработал на Go (gRPC, PostgreSQL, Kafka) сервис сборки заказа для конфигуратора ПК: он агрегирует комплектующие, формирует инвойс и передаёт заказ в сервис оформления; держит до 1\,200 RPS в пике распродаж при p99 180 мс.
  \item Сократил p95 формирования инвойса с 420 до 150 мс (−64\%), заменив последовательные вызовы сервисов цен, остатков и совместимости на параллельные через errgroup с Redis-кэшем.
  \item Уменьшил долю потерянных и задублированных заказов при передаче в оформление с 0,8\% до 0,05\%, внедрив transactional outbox, ключи идемпотентности и ретраи с backoff.
\end{itemize}

% Несколько ролей в одной организации: каждая должность со своими датами.
% \cventry{Организация}{Месяц 20XX -- Месяц 20XX}{}{Город, Страна}
% \cvrole{Старшая должность}{Месяц 20XX -- Месяц 20XX}
% \begin{itemize}
%   \item Пункт.
% \end{itemize}
% \cvrole{Должность}{Месяц 20XX -- Месяц 20XX}
% \begin{itemize}
%   \item Пункт.
% \end{itemize}

% Контракт: явно помечай как контракт.
% \cventry{Клиент или Самозанятость}{Месяц 20XX -- Месяц 20XX}{Должность (контракт)}{Удалённо}

% Перерыв в карьере: по желанию, без медицинских и семейных подробностей.
% \cventry{Перерыв в карьере}{Месяц 20XX -- Месяц 20XX}{Учёба, уход за близкими, переезд}{}

\section*{Сторонние проекты}
\cventry{ReadBeat}{Июнь 2026 -- н.\,в.}{CTO}{Удалённо}
\cvcontext{Медиапортал битбокс-сообщества: афиша баттлов, фестивалей и джемов.}
\begin{itemize}
  \item Спроектировал бэкенд на Go: GraphQL API (gqlgen), REST-загрузка изображений с документацией OpenAPI, PostgreSQL (pgx, golang-migrate) и сессии администраторов в Redis.
  \item Закрепил слои Clean Architecture проверками arch-go и golangci-lint (gosec, revive) в CI.
  \item Построил пайплайн GitHub Actions: проверка актуальности кодогенерации, интеграционные тесты с race detector на PostgreSQL и Redis, деплой distroless Docker-образов за Nginx по SSH.
\end{itemize}

\cventry{Big Dope Label}{Март 2026 -- н.\,в.}{CTO}{Удалённо}
\cvcontext{Первый битбокс-лейбл в России --- \href{https://bigdope.pro}{bigdope.pro}.}
\begin{itemize}
  \item Настроил продакшен-хостинг: Nginx в Docker с TLS от Let's Encrypt и автоматическим продлением сертификатов.
  \item Собрал аналитику трафика и управление логами в отдельных контейнерах: GoAccess, logrotate и просмотрщик логов под паролем.
  \item Автоматизировал CI/CD на GitHub Actions: линтеры и тесты на каждый push, деплой через Docker Compose по SSH.
\end{itemize}

\section*{Образование}
% Квалификация, вуз, город, год окончания. Для незавершённой программы —
% ожидаемая дата; «Обучение по программе ...», если не окончил.
% Средний балл — только если высокий и свежий, со шкалой (4.8/5.0).
\cventry{Бакалавр, Математика и компьютерные науки}{2022}{Дальневосточный федеральный университет}{Владивосток, Россия}

% \section*{Сертификаты}
% Два-три релевантных: официальное название, кто выдал, дата выдачи/окончания,
% ID или ссылка. «Пройден курс», если экзамена не было.
% \textbf{Название сертификата}, Организация, Месяц 20XX\hfill ID: XXXX

\section*{Языки}
% Шкала CEFR (A1-C2) или «родной». Не завышай уровень.
Английский (B2)\cvsep Русский (родной)

% \section*{Интересы}
% По желанию и коротко: только конкретные, весомые достижения.
% Пробежал X и организовал Y для N человек.

\end{document}
